\emph{Setting.} The bivectors of the block $e_3, \dots, e_8$ span $\so(6)$; those commuting with $J$ span $\uu(3)$ (nine dimensions), and those orthogonal to $J$ for $\langle X\widetilde{Y}\rangle_0$ span the $\su(3)$ of \claimref{synthesis.field-level-brst-differential-of-the-derived-gauge-fields-is-nilpotent}(i), which proves it simple of dimension $8$, isomorphic to the colour summand of Theorem \claimref{standard-model.gauge-algebra-minimal-massive-composite}. On vectors $j(v) = v\cdot J$ sends $e_3 \mapsto e_4 \mapsto -e_3$ (and likewise on $e_5, e_6$ and $e_7, e_8$) and is the complex structure; $\Omega_R$ and $\Omega_I$ are the real and imaginary parts of the complex volume form on the frame $e_3, e_5, e_7$, whose $j$-partners are $e_4, e_6, e_8$. The check multiplies even elements of the block with a blade table built from GA's blade product and compares it with the Clifford product on all $32^2$ basis pairs.

(i) The $P_k$ commute with each other and with $I$, and $(IP_k)^2 = I^2P_k^2 = 1$, so the $\tfrac12(1 + IP_k)$ are commuting idempotents and $E^2 = E$; $\widetilde{I} = -I$ and $\widetilde{P_k} = -P_k$, so $\widetilde{IP_k} = P_kI = IP_k$ and $\widetilde{E} = E$. Expanding the product, the scalar terms are $1$ and $I^3P_{34}P_{56}P_{78} = I^4 = 1$, so $\langle E\rangle_0 = \tfrac14$. $P_k(1 + IP_k) = P_k - I = -I(1 + IP_k)$ gives $P_kE = -IE$, hence $(P_j - P_k)E = 0$ and $JE = -3IE \neq 0$. \emph{Computed:} on the fifteen bivectors the kernel of $B \mapsto BE$ is eight-dimensional and spans the same space as the eight $\su(3)$ bivectors; every $\su(3)$ bivector commutes with $J$, $\Omega_R$ and $\Omega_I$, while $[J, \Omega_R] = -6\Omega_I$ and $[J, \Omega_I] = 6\Omega_R$.

\emph{Why the vector condition.} In six dimensions $R\widetilde{R} = 1$ does not make an even element a rotor: with $\cos\tfrac\varphi2 = \tfrac35$, $R = \exp(-\tfrac32\varphi I)\exp(-\tfrac\varphi2J)$ is rational and even, $R\widetilde{R} = 1$ and $RE = E$ (since $\exp(-\tfrac\varphi2J)E = \exp(\tfrac32\varphi I)E$), but $Re_3\widetilde{R}$ has parts of grades $1$ and $5$ (computed). Elements with $R\widetilde{R} = 1$ fixing $E$ form a $U(3)$ on which $R \mapsto f_R$ is not a rotation; requiring $Rv\widetilde{R}$ to be a vector for every $v$ removes them. \emph{Colour rotors commute with $J$.} If $RE = E$ then $\widetilde{R}E = E$, and $f_R(B)E = RB\widetilde{R}E = R(BE)$ for every bivector $B$. For $B = J$ this is $R(-3IE) = -3IE = JE$, so $f_R(J) - J$ lies in the kernel, $\su(3)$, which is orthogonal to $J$. Writing $f_R(J) = J + S$, $\langle f_R(J)\widetilde{f_R(J)}\rangle_0 = \langle J\widetilde{J}\rangle_0$ gives $\langle S\widetilde{S}\rangle_0 = 0$, so $S = 0$ (the form is positive definite on bivectors of $\Cl(9,0)$).

\emph{Determinant.} A rotation $f$ commuting with $j$ is complex linear and orthogonal, so unitary on $(\R^6, j) \cong \C^3$. $U(3)$ is the product of $SU(3)$ and the scalars $\exp(\varphi j) = \cos\varphi + \sin\varphi\,j$, and the scalars are the $J$-rotations with rotors $\pm\exp(-\varphi J/2)$. On trivectors $\su(3)$ acts trivially on $\Omega_R, \Omega_I$ and $J$ turns them by the brackets above, so $\exp(-\varphi J/2)$ sends $\Omega_R$ to $\cos3\varphi\,\Omega_R + \sin3\varphi\,\Omega_I$, and every $f = \exp(\varphi j)g$ with $g \in SU(3)$ (generated by $\exp\su(3)$) does the same, its determinant being $\exp(3\varphi j)$; so $f(\Omega_R) = \Omega_R$ exactly when $3\varphi \in 2\pi\Z$, that is $\det f = 1$. \emph{Computed:} the quarter turn $j$ itself ($\varphi = \tfrac\pi2$, determinant $-j$) sends $\Omega_R \mapsto -\Omega_I$ and $\Omega_I \mapsto \Omega_R$. At rotor level, a rotor of $f = \exp(\varphi j)g$ is $\pm\exp(-\varphi J/2)R_g$ with $R_g$ a product of exponentials of $\su(3)$ bivectors, so $R_gE = E$; and $\exp(-\varphi J/2)E = \exp(\tfrac32\varphi I)E$ since $JE = -3IE$. Hence $RE = E$ forces $\exp(\tfrac32\varphi I) = \pm1$, $\varphi \in \tfrac{2\pi}{3}\Z$, $\det f = 1$: a colour rotor covers an element of $SU(3)$. Conversely every element of $SU(3)$ is a product of exponentials of $\su(3)$ (it is connected), whose rotors fix $E$; and of the two rotors $\pm R$ of a rotation only one fixes $E$. So $R \mapsto f_R$ is an isomorphism of the colour rotors onto $SU(3)$.

(ii) Both sides are class functions. Every element of $SU(3)$ is conjugate in $SU(3)$ to a diagonal one (an orthonormal eigenbasis, with one vector rescaled by a phase so the conjugating matrix has determinant $1$), whose colour rotor is $T_\theta = \prod_k(\cos\tfrac{\theta_k}{2} - \sin\tfrac{\theta_k}{2}P_k)$ with $\theta_1 + \theta_2 + \theta_3 = 0$: by $P_kE = -IE$, $T_\theta E = \exp(\tfrac12(\theta_1 + \theta_2 + \theta_3)I)E$, which is $E$ for this choice. Expanding, $\langle T_\theta\rangle_0 = c_1c_2c_3$ and $\langle T_\theta\rangle_6 = -s_1s_2s_3I$ ($c_k, s_k$ the half-angle cosines and sines; products of two $P_k$ have grade four). With $\sum\theta_k = 0$ the product formulas give $4c_1c_2c_3 = 1 + \sum_k\cos\theta_k$ and $-4s_1s_2s_3 = \sum_k\sin\theta_k$, and $f_{T_\theta}$ rotates each plane $P_k$ by $\theta_k$, so it is the scalar $\exp(\theta_kj)$ on the $k$-th plane. So $\chi = 4(\langle R\rangle_0 + \langle R\rangle_6) - 1$ is the trace of the triplet, $\sum_k(\cos\theta_k + I\sin\theta_k)$, written with $I$ standing for $j$. For a $j$-linear $f$ the trace, so written, is $\tfrac12\tr_{\R}f - \tfrac12\tr_{\R}(j\circ f)\,I$, which gives the two vector formulas. \emph{Computed:} on all $96$ elements of the group below and on $36$ conjugates of three rational torus rotors. So $w = \tfrac13\mathrm{Re}\Tr$, the weight is Wilson's $\exp(-\beta\sum_p(1 - \tfrac1N\mathrm{Re}\Tr U_p))$ up to a constant with $N = 3$, and $\beta = 6/g^{2}$ (Fact \claimref{synthesis.imported-wilson-lattice-gauge-theory-and-strong-coupling-confinement}(i), $\beta = 2N/g^{2}$); $1 - w = \tfrac43(1 - \langle U\rangle_0)$ is immediate.

(iii) This lemma supplies the moments the later clauses use. \emph{Averages are projectors.} For a compact group with normalised Haar measure and an orthogonal representation $\rho$, $P = \int\rho(R)\,dR$ satisfies $\rho(S)P = P$ (left invariance), $Pv = v$ for every invariant $v$, and $P^{T} = \int\rho(R^{-1})\,dR = P$ (inversion invariance), so $P$ is the orthogonal projector onto the invariants; the same holds for the uniform average over a finite group. \emph{The group.} $G$ is generated by the colour rotors of the quarter turns acting as $(j, -j, 1)$ and $(1, j, -j)$ on the three planes ($\tfrac12(1 - P_{34})(1 + P_{56})$ and $\tfrac12(1 - P_{56})(1 + P_{78})$), the cyclic permutation of the three complex coordinates and minus a transposition (products of $\tfrac12(e_3 - e_5)(e_4 - e_6)$, $\tfrac12(e_5 - e_7)(e_6 - e_8)$ and $P_{78}$), each with the sign that fixes $E$. \emph{Computed:} $G$ has $96$ elements, all rational, all fixing $E$, $-1 \notin G$; their vector actions are $J$-commuting signed permutation matrices (this is $\Delta(96)$). By (i) $G$ is a subgroup of the colour rotors, so every $SU(3)$-invariant tensor is $G$-invariant. \emph{Computed:} the dimension of the $G$-invariants of $f^{\otimes k}$, the average of $(\tr f)^{k}$, is $0, 2, 2$ for $k = 1, 2, 3$. In degree $2$ the tensors $\delta$ and $j$ are fixed by every colour rotor ($f$ orthogonal and commuting with $j$), in degree $3$ the totally antisymmetric tensors of $\Omega_R$ and $\Omega_I$ are (by (i)); each pair is independent. So in each degree the $SU(3)$-invariants and the $G$-invariants coincide, the two projectors are equal, and the Haar moments are the $G$-averages, which the check computes entry by entry: $0$ in degree $1$, $\tfrac16(\delta\delta + jj)$ in degree $2$ (norms $|\delta|^2 = |j|^2 = 6$), and $\tfrac1{24}(\Omega_R\Omega_R + \Omega_I\Omega_I)$ in degree $3$ (tensor norm $24$), which in wedge form is the displayed formula with $\langle\Omega_R\widetilde{\Omega_R}\rangle_0 = 4$, checked on all twenty basis trivectors. For left multiplication on the even subalgebra (orthogonal for $\langle X\widetilde{Y}\rangle_0$), $\tr L_R = 32\langle R\rangle_0$ averages to $8$ over $G$, and the right ideal $E\Cl^{+}$, fixed by every colour rotor, has dimension $8$; so the projectors agree there too and $\int R\,dR = \int L_R(1)\,dR$ is the $G$-average of the rotors, which is $E$. \emph{Characters.} For $J$-commuting $X$, $Y$ (so $jYj^{T} = Y$) the two terms of the degree-two moment give $\int\tr_{\R}(Xf)\tr_{\R}(f^{T}Y) = \tfrac16(\tr_{\R}(XY) + \tr_{\R}(jYj^{T}X)) = \tfrac13\tr_{\R}(XY)$, hence $\int w(XR)w(\widetilde{R}Y) = \tfrac1{18}w(XY)$ with $w = \tfrac16\tr_{\R}$; read on $\C^3$, the real invariants $\delta$ and $j$ are the real and imaginary parts of the one invariant of the triplet times the antitriplet, and the triplet times itself has none, so $\int\chi(XR)\chi(\widetilde{R}Y) = \tfrac13\chi(XY)$ and $\int\chi(XR)\chi(YR) = 0$; the degree-three moment gives $\int\chi^3 = 1$ (one invariant, the volume form) and $\int\chi^2\bar\chi = 0$. \emph{Computed:} the five single-link averages, and the three gluing identities for all $96^2$ pairs $(X, Y)$ in $G$; both sides are bilinear in the complex matrices of $X$ and $Y$, and the vector actions of $G$ span the eighteen-dimensional space of $J$-commuting matrices (Burnside), so the check on $G \times G$ is complete. \emph{Contrasts.} The cubic moment averages to zero over the $J$-rotations, which turn the pair $(\Omega_R, \Omega_I)$ by the angle $3\varphi$ (checked on the quarter turns), so over $U(3)$ it vanishes; for the weak factor every odd moment vanishes ($-1 \in SU(2)$). Nothing below uses moments of degree $4$.

(iv) $z$ is $-T_\theta$ for $\theta = (\tfrac{2\pi}{3}, \tfrac{2\pi}{3}, \tfrac{2\pi}{3})$, that is $T_\theta$ for $(\tfrac{2\pi}{3}, \tfrac{2\pi}{3}, -\tfrac{4\pi}{3})$, so it is a colour rotor; its vector action multiplies by the scalar $\omega = -\tfrac12 + \tfrac{\sqrt3}{2}j$, so $z$ is central, $z^3 = 1$, and $w(z) = \mathrm{Re}\,\omega = -\tfrac12$. \emph{Computed} in $\Q(\sqrt3)$: $zE = E$, $z^3 = 1 \neq z^2$, $\langle z\rangle_0 = -\tfrac18$, $z$ commutes with the generators of $G$ and with all eight $\su(3)$ bivectors, and $f_z(e_k) = -\tfrac12e_k + \tfrac{\sqrt3}{2}e_{k+1}$ for $k = 3, 5, 7$. By left invariance under $R \mapsto zR$ the displayed moment equals $\omega^{a - b}$ times itself, so it vanishes unless $a \equiv b \pmod 3$. On the lattice, expand $\exp(\beta w) = \exp(\tfrac\beta6\chi)\exp(\tfrac\beta6\bar\chi)$, so each term is a product over plaquettes of monomials $\chi(U_p)^{a_p}\bar\chi(U_p)^{b_p}$ of order $a_p + b_p$, times the observable's monomials; give $p$ the charge $n_p = a_p - b_p \bmod 3$. The change of variables $U_\ell \mapsto zU_\ell$ on one link preserves the Haar measure and multiplies such a term by $\omega$ to the net charge through $\ell$ ($z$ central, so it passes to the front of every plaquette product containing $\ell$, with $\bar\omega$ for the reversed orientation). So a term survives only if $n = \sum_pn_pp$ is a $\Z_3$ two-chain with $\partial n$ equal to minus the observable's charge chain, and its order is at least the number of plaquettes with $n_p \neq 0$.

(v) \emph{No lower order.} The observable $w(U_C) = \tfrac16(\chi + \bar\chi)(U_C)$ has charge $\pm1$ on $C$, so $\partial n = \mp C$. The projection onto the plane of $C$ along the two other axes is a chain map; it sends plaquettes parallel to the plane to plaquettes and the others to degenerate cells, so it maps $n$ to a $\Z_3$ two-chain of the plane with boundary $\mp C$. The plane has no nonzero finite two-cycle, so this chain is $\mp$ the sum of the $A$ plaquettes inside $C$, and over each of them $n$ has a plaquette with $n_p \ne 0$: every surviving term has order at least $A$. At order $A$ exactly there is one such plaquette over each interior position, all at first order and none elsewhere; a plaquette next to $C$ must lie in the plane of $C$ to produce its edge, and two plaquettes over adjacent positions must lie at the same transverse position to cancel their common edge, so by connectedness the plaquettes are those of the rectangle. The normalisation is $1 + O(\beta^2)$, since $\int w = 0$. So the coefficient of $\beta^{A}$ is $\int w(U_C)\prod_{p \subset C}w(U_p)$, the only term of that order. \emph{Its value.} As for the weak factor: fix the tree of all horizontal links and the left vertical links to $1$, change variables row by row to the independent Haar rotors $P_p = V_{x+1,y}\widetilde{V_{x,y}}$, and integrate one $P_p$ at a time with the gluing rule of (iii) at $Y = 1$ (using $w(\widetilde{P}) = w(P)$); each step gives $\tfrac1{18}$, so the value is $(\tfrac1{18})^{A}$ and the leading coefficient is $(\beta/18)^{A}$. \emph{Computed:} $\tfrac1{18}$ and $\tfrac1{324}$ for the $1 \times 1$ and $2 \times 1$ loops by summing over $G$ in the tree gauge; independently, $18^{-1}$, $18^{-3}$, $18^{-4}$ for $1 \times 1$, $3 \times 1$, $2 \times 2$ by contracting the degree-two moment over every link of the disc (each link lies on exactly two of the faces, the loop counted as a face). Gauge invariance of closed loops and the failure for an open line are checked on a $2 \times 2$ patch with site rotors in $G$. \emph{$N$-ality zero.} Three fundamental lines along the same closed path joined by the volume form at the base point are, by complex multilinearity, $\tfrac14\langle f_{U_C}(\Omega_R)\widetilde{\Omega_R}\rangle_0$, which is $1$ for every configuration by (i): no suppression at any order, in contrast with the bound above, which needs nonzero charge.

(vi) The box has periodic time and free spatial boundary (a periodic spatial extent would let a chain wind the short way round). \emph{Exact zeros.} Multiplying every temporal link of one time slice by $z$ preserves the measure and every plaquette (a temporal plaquette contains one such link and the inverse of another, and $z$ is central), and multiplies $P(x)$ by $\omega$; so $\langle P(x)\rangle = \omega\langle P(x)\rangle = 0$ and $\langle P(x)P(y)\rangle = \omega^2\langle P(x)P(y)\rangle = 0$ at every $\beta$, while $P\bar P$ and $PPP$ are invariant. \emph{No lower order.} In each time slab the temporal plaquettes are the edges of the spatial lattice, and the charge condition at the temporal links of the slab says that their charges form a $\Z_3$ one-chain with boundary $\pm(x - y)$, or $\pm(a + b + c)$ for three lines at $a, b, c$. For $y = x + re_i$, projecting on the $i$-th axis shows that such a chain has at least $r$ edges, with equality only for the straight segment; for $a, b, c = x + e_1, x + e_2, x + e_3$, which are pairwise not adjacent, at least three edges are needed, each containing one of them, and their other ends must carry boundary $0 \bmod 3$, which forces a common end, the only common neighbour $x$ (charge $3 \equiv 0$: the baryonic junction). So the order is at least $N_tr$, respectively $3N_t$, and at that order there are no spatial plaquettes, whence the spatial links at each time force the same strip (respectively star) in every slab. \emph{Values.} In each strip integrate the spatial links one by one with $\int\chi(XS)\chi(\widetilde{S}Y) = \tfrac13\chi(XY)$, and the last with $\int\chi(XSY\widetilde{S}) = \tfrac13\chi(X)\chi(Y)$ (from the degree-two moment): a strip of $N_t$ plaquettes of width one, each carrying $\tfrac\beta6\chi$ with the orientation that pairs with the lines, becomes $(\tfrac\beta{18})^{N_t}\bar\chi(L_u)\chi(L_v)$ for its two edges $u$, $v$, with $L$ the line products, which are independent Haar rotors. Gluing the $r$ columns with $\int\chi\bar\chi = 1$ on the inner lines and pairing with $P(x)\overline{P(y)} = \tfrac19\chi(L_x)\bar\chi(L_y)$ gives $\tfrac19(\beta/18)^{N_tr}$; for the junction the three sheets give $(\tfrac\beta{18})^{3N_t}\chi(L_x)^3\prod_k\bar\chi(L_{x+e_k})$, and $\int\chi^3 = 1$ (the baryonic moment) with three factors $\int\chi\bar\chi = 1$ gives $\tfrac1{27}(\beta/18)^{3N_t}$. \emph{Computed} by sums over $G$ in a tree gauge: $\tfrac1{162}$ for $N_t = r = 1$, $\tfrac19\cdot18^{-2}$ for $N_t = 2$, $r = 1$, zero for $P(x)P(y)$ at the same order, each sheet of the junction reducing to $\chi(U_x)/54$ for every $U_x \in G$, and $\tfrac1{27}\cdot18^{-3}$ for the junction at $N_t = 1$.

(vii) Write $\exp(\beta w(U)) = c_0(1 + h(U))$ with $c_0 = \int\exp(\beta w)$, so $\int h = 0$ and $h = \beta w + O(\beta^2)$. In a box with free boundary expand $\prod_p(1 + h_p)$ over sets $S$ of plaquettes; the term of $S$, with the observables $O_0 = w(U_p)$, $O_t = w(U_{p_t})$ inserted or not, vanishes if some link lies in exactly one member of $S \cup \{p, p_t\}$ (or of $S$ with the observables absent), by integrating that link ($\int h_q(AU_\ell B)\,dU_\ell = \int h = 0$ and $\int w = 0$). Put a parameter $s_\tau$ on every $h_q$ with $q$ a temporal plaquette of slab $\tau$, $0 \le \tau < t$. The connected correlator $C = N/Z - N_0N_t/Z^2$ is a formal power series in $\beta$ with coefficients polynomial in the $s_\tau$. (a) At $s_\tau = 0$ no remaining factor contains a temporal link of slab $\tau$ and the links split into those at times $\le \tau$ and those at times $\ge \tau + 1$, so the measure factorises with $O_0$ below and $O_t$ above, and $C = 0$; so every monomial of $C$ contains every $s_\tau$. (b) In a nonvanishing term the slab-$\tau$ plaquettes of $S$ form, as edges of the spatial lattice, a graph with no vertex of degree one (their temporal links meet no observable), so if nonempty it contains a cycle, of length at least $4$ in the cubic lattice; every term of $N$, $Z$, $N_0$ and $N_t$, and hence every product in the expansion of $C$, has $s_\tau$-degree $0$ or at least $4$. With (a), every monomial of $C$ has degree at least $4$ in each $s_\tau$, and since $h = O(\beta)$ its order is at least $4t$. At order $4t$ every factor is $\beta w$ and there are exactly four per slab and none elsewhere. A term of $N_0$ or $N_t$ with slab plaquettes would leave the spatial links at the top (or bottom) of its last slab in one plaquette only, so $N_0N_t/Z^2$ contributes nothing, and $Z - 1 = O(\beta^2)$ adds order; in $N$ the four plaquettes of each slab are a unit square, and the spatial links at each time force all squares to be the boundary of $p$. So the coefficient of $\beta^{4t}$ is $\int O_0O_t\prod_{\mathrm{tube}}w$. \emph{Its value.} By (iv) only the two consistent orientations survive; for each, integrating the four temporal links of a slab with the gluing rules gives $3^{-4}\chi(S_\tau)\bar\chi(S_{\tau+1})$ for the squares $S_\tau$ at the two ends, and $\int\chi\bar\chi = 1$ joins consecutive slabs and the ends: $3^{-4t}$, so the coefficient is $2\cdot6^{-(4t+2)}\cdot3^{-4t} = \tfrac1{18}(\tfrac1{18})^{4t}$, and $-\tfrac1t\ln$ of the leading term tends to $4\ln(18/\beta)$. \emph{Computed:} $18^{-5}$ and $18^{-9}$ for $t = 1, 2$ by contracting the degree-two moment over every link of the closed tube.

\emph{What is proved and what is not.} Proved: (i)-(vii) as stated, each at the stated order on a finite lattice; the Haar moments of degree at most three without importing them; the leading terms and the absence of lower-order terms. Not proved: the remainders' uniformity in the separation (for (vii), a gap of size $4\ln(18/\beta) + O(\beta)$ at small $\beta$ in infinite volume), which is expected to follow from the convergent strong-coupling expansion of Osterwalder and Seiler (cited from memory, not imported; their area law and analyticity are Fact \claimref{synthesis.imported-wilson-lattice-gauge-theory-and-strong-coupling-confinement}(ii)). \emph{Argued, not proved here:} the centre bound of (iv) does not apply to loops of $N$-ality zero (the cases of (v) and (vi) are examples), and the adjoint loop is screened at large size, its leading terms growing with the perimeter. The same leading structure holds for compact $U(1)$ (the weak-factor result computes $(\beta/2)^{A}$), and four-dimensional compact $U(1)$ deconfines at weak coupling (Fact (iv)); so these strong-coupling statements are not evidence of confinement in the continuum. The continuum limit needs $\beta \to \infty$, where the expansion does not converge (the leading rate $4\ln(18/\beta)$ is not small there), and asymptotic freedom (the store's one-loop running coefficients from spacetime-algebra spin traces); the Yang-Mills existence and mass-gap problem is a statement about that limit and is not addressed. Transport-absent.
